\subsection{All Event Types: Network Physics Without Training}
\label{sec:allevents}

Because the bank cannot start on the competition record
(Section~\ref{sec:integrity-results}), we built a second attributor from
the provided network data alone, with no trained parameter, for every
event type of the guide. The
network is the admittance matrix of the provided power-flow case, with
loads as constant impedances and generators as transient reactances
behind constant voltages~\cite{pai1989energy}. The integrity layer runs
first; unhealthy PMUs are deleted from the fits
(Proposition~\ref{prop:integrity}), not imputed. The physical layer uses
the two time scales of the network equations. \emph{Algebraic jump:}
the one-report jump of the PMU voltages and branch currents is fitted by
an injection at each of the 39 buses or the removal of each of the 46
branches, removed with its full pre-event current at both terminals; a
sag below 0.8\,pu is a fault, located by the profile $|Z_{jf}/Z_{ff}|$.
\emph{Quasi-steady state:} an episode that does not return within 15\,s is
fitted 15--25\,s after onset by the linearized AC power flow, with
generators sharing the imbalance in proportion to dispatch. This
separates load changes, generation changes and generator trips, under
plausibility limits (a load cannot change by more than its base value,
nor a generator by more than its dispatch). A reversed jump, or a return
to the pre-episode state, closes an episode. Attributor~I detects
episodes from the change of the angle differences. Attributor~II detects
one-report jumps after subtracting the local trend of the neighboring
report changes, which separates switching from the smooth
electromechanical response. Decisions are taken 0.5\,s (switching) or
25\,s (persistent) after onset and written back to it, as in the guide's
offline scoring; each record takes seconds.

\emph{Protocol.} No label was used to fit a parameter, but the
competition record was not blind. Attributor~I took \AERawIterOne{}
logged iterations on it. Attributor~II was designed on the training and
calibration records of a simulated cohort and checked on the competition
record in two further iterations. Every iteration is stored with its code hash
and score; the competition record is thus a validation set.
The confirmatory test is the stored test split of that cohort, generated
before either attributor existed. Each attributor was fixed and run
once, and its predictions were hashed before the schedule was read.

\emph{Results.} On the competition record (Table~\ref{tab:allevents-raw})
both attributors detect, classify and locate all 12 labeled events of that one record. The
submitted system located \AERawSubTopone{} of them. The generation event
is the trip of the machine at bus~30 (fitted scale 1.02, no amplitude
fit); the load at bus~7 is separated from generation by its
reactive-to-active ratio. Six of the eight false-alarm episodes of
Attributor~II are the impossible measurements of
Section~\ref{sec:integrity-results}, a seventh follows the collapse at
PMU~29, and the eighth is an unlabeled change at 3580\,s seen at every
PMU. Labeled boundaries that extend beyond the
observable change keep macro-F1 at \AERawTwoMacro{}.
\begin{table}[t]\centering\footnotesize
\caption{All event types on the competition record (12 labeled events: one fault, one line, one generation and one load change, four dropouts, three corruptions). Report-level scores follow the competition guide; false alarms are episodes outside every label.}
\label{tab:allevents-raw}\setlength{\tabcolsep}{3.2pt}
\begin{tabular}{lccccc}\toprule
Method & Class & Top-1 & F1 & FA/min & Macro-F1 \\
\midrule
Submitted system & 11/12 & 7/12 & 0.192 & 0.18 & 0.21 \\
Attributor I & 12/12 & 12/12 & 0.981 & 0.10 & 0.90 \\
Attributor II & 12/12 & 12/12 & 0.979 & 0.10 & 0.88 \\
\bottomrule\end{tabular}\end{table}

\begin{table}[t]\centering\footnotesize
\caption{Stored simulated cohort, evaluated once: 72 load and line pulses (2--5\,s, six per record) in 12 records, and 12 normal records (2\,h). Top-1 counts events with the right family and exact source. The dynamic bank uses the plant model and an 18-source catalog; the open search ranks all 39 buses and 46 branches.}
\label{tab:allevents-sim}\setlength{\tabcolsep}{2.6pt}
\begin{tabular}{lcccccc}\toprule
Method & Detected & Top-1 & Lines & Loads & Loads & Normal \\
 & & & Top-1 & detected & Top-1 & alarms \\
\midrule
Attributor I, open & 13/72 & 0/72 & 0/36 & 4/36 & 0/36 & 0 \\
Attributor II, open & 48/72 & 28/72 & 24/36 & 12/36 & 4/36 & 0 \\
Attributor II, catalog & 48/72 & 40/72 & 32/36 & 12/36 & 8/36 & 0 \\
Dynamic bank, catalog & 72/72 & 72/72 & 36/36 & 36/36 & 36/36 & 0 \\
\bottomrule\end{tabular}\end{table}


The stored cohort has 72 load and line pulses of 2--5\,s, six per record
with gaps of 3--8\,s (Table~\ref{tab:allevents-sim}). Attributor~I fails
there: its quasi-steady window opens after the pulses end. Attributor~II
detects all \AESimLineDet{} line pulses and locates \AESimLineTopone{}
in an open search over 85 candidates, or \AESimCatLineTopone{} within
the bank's 18-source catalog. Every remaining line error involves one of
three branches whose outage, seen from the eight PMUs, is nearly
collinear with an injection at an adjacent bus. It detects only
\AESimLoadDet{} of the 36 load pulses, because a change of 3--14\% of one
load moves the one-report phasors by less than the startup noise. Longer
windows would recover them, but there the response is no longer
algebraic. The dynamic bank models that response and attributes all 72
pulses, at \AEWallBankH\,h of computation against \AEWallTwo\,s. No
method alarmed in the normal records.

Network physics alone thus attributes every event type that is large
enough for its algebraic jump or long enough to settle; short, small
pulses need the dynamic model.
